\subsection{Impact of Input Characteristics}
\label{sec:topology_impact}

% Notes for relative time measurement:
% Times are normalized independently within each graph size category (1K, 10K, 100K, 2M) by dividing by the mean execution time of that scale. The red dashed line ($y = 1.0$) denotes the average computational cost for a graph of that size, highlighting which structural traits inflate or reduce execution overhead.
% For example in subplot (a), for 1K S instances, the mean of all the 6 blue bars is taken and for each bar the time there is divided by this mean. This is done for all S, L, XXL, XXXL seperately

% 

BP-MDS is evaluated on diverse input characteristics (Fig.~\ref{fig:characteristics}) generated using the CVRPLIB input generator. Execution time decreases as the average route size increases, since longer routes reduce the overhead of repeatedly switching between buckets in a parallel setting (Fig.~\hyperref[fig:characteristics]{\ref*{fig:characteristics}(a)}). For the 2M-node instances, increasing the average route size from \emph{Very Short} to \emph{Ultra Long} reduces execution time from 204.04\,s to 151.68\,s. Cornered-depot instances generally require higher execution time than centered or randomly positioned depots (Fig.~\hyperref[fig:characteristics]{\ref*{fig:characteristics}(b)}). This is because cornered depots produce unbalanced bucket sizes, increasing the workload of the largest bucket. On 1K (S) graphs, the cornered depot incurs a $2.47\times$ higher execution time than the centered depot. This behavior is evident in Fig.~\hyperref[fig:bucket_analysis]{\ref*{fig:bucket_analysis}(a,b)}, where \texttt{Brussels2} and \texttt{Flanders2} (cornered depots) exhibit significantly larger buckets than \texttt{Brussels1} and \texttt{Flanders1} (centered depots) (spatial differences are visualized in Fig.~\ref{fig:BKS_Brussels}). Execution time strictly increases with the maximum bucket size (Fig.~\hyperref[fig:bucket_analysis]{\ref*{fig:bucket_analysis}(c)}), confirming that the largest bucket dominates the overall runtime. \REM{In contrast, customer positioning (Fig.~\hyperref[fig:characteristics]{\ref*{fig:characteristics}(c)}) and demand distribution (Fig.~\hyperref[fig:characteristics]{\ref*{fig:characteristics}(d)}) have minimal impact on execution time.}

% {\color{magenta} Demand Distribution among customers. Type 1 assigns a unitary demand of exactly $1$ to all customers, while Types 2 through 5 sample uniformly from fixed intervals: $[1, 10]$ (small, large variance), $[5, 10]$ (small, small variance), $[1, 100]$ (large, large variance), and $[50, 100]$ (large, small variance). For Quadrant, Customers in "even quadrants" (bottom-left where $x,y < 5000$, or top-right where $x,y \ge 5000$) get large demands sampled uniformly from $[51, 100]$.  Customers in the other two quadrants get smaller demands sampled from $[1, 50]$. A calculated proportion of customers ($\frac{n}{r} \times 1.5$, where $r$ is target route size) receive large demands uniformly sampled from $[50, 100]$.  All remaining customers receive small demands sampled from $[1, 10]$. 
% Average Route Size (avgRouteSize) controls the expected number of customers served per vehicle by uniformly sampling a target continuous value $r$(route size) from one of six category intervals: Very short $[3, 5]$, Short $[5, 8]$, Medium $[8, 12]$, Long $[12, 16]$, Very long $[16, 25]$, and Ultra long $[25, 50]$. 
% Depot Positioning establishes the starting hub location on a $10{,}000 \times 10{,}000$ 2D Euclidean grid according to three simple placement strategies. Type 1 (Random) places the depot uniformly at random anywhere within the $[0, 10000]$ coordinate range. Type 2 (Centered) fixes the depot exactly at the midpoint of the grid at coordinates $(5000, 5000)$. Type 3 (Cornered) fixes the depot at the bottom-left origin point $(0, 0)$.  
% Customer Positioning (custPos) distributes the $n$ customer coordinates across the same $10{,}000 \times 10{,}000$ Euclidean grid using three distinct spatial patterns. Type 1 (Random) generates all customer coordinates uniformly at random across the entire grid. Type 2 (Clustered) generates between 2 to 6 random seed centers and places customers around them using an Accept-Reject method. Type 3 (Random-Clustered) is a 50/50 hybrid model where half of the customers ($n/2$) are generated uniformly at random and the remaining half are generated around seed centers using the exponential decay clustering method.}

%  The below two paras are related to: Impact of graph topology on alpha. This is removed because of lack of space. 
% Fig.~\hyperref[fig:characteristics]{\ref*{fig:characteristics}(a)} shows that $\alpha$ increases with average route size, since larger routes require a wider angular sweep to accumulate sufficient customers for feasible vehicle capacity. Fig.~\hyperref[fig:characteristics]{\ref*{fig:characteristics}(b)} shows that a centered depot yields higher $\alpha$ than cornered or random placements, as central positioning spreads customers over a full $360^\circ$ range, requiring wider partitions to capture enough nodes per bucket. Fig.~\hyperref[fig:characteristics]{\ref*{fig:characteristics}(c)} shows that clustered customer distributions also increase $\alpha$, as overly narrow partitions fragment spatial clusters and break short-distance route coherence.

% Across all cases, graph scale dominates these effects: as the number of nodes increases, sensitivity to topology diminishes and $\alpha$ consistently converges to smaller angular values. This is because higher node density reduces the need for wide partitions, allowing efficient routing within narrow angular regions. Overall, BP-MDS adapts to topology at small scales but becomes scale-dominated, where dense graphs enforce uniform, highly fine-grained partitioning without loss of routing quality.
